\chapter{Symmetry Actions and Represented Subspaces}
\label{ch:symmetry-actions-represented-subspaces}
\index{symmetry!represented action|textbf}
\index{representation!finite group}

\chapterstatus{finite-dimensional mathematical foundations, not a report of
implemented symmetry software or electronic symmetry extraction. The sector
results below concern ordinary complex-linear representations of finite groups.
They establish no silicon symmetry-sector attribution and claim no new
representation-theory theorem. The sector results remain manuscript proofs;
the separate coordinate-algebra mechanization note below has a strictly
narrower scope.}

\section{Three distinct symmetry statements}

A geometric operation preserving a crystal motif is not yet an action on its
electronic states. An electronic action preserving the inner product is not yet
a symmetry of a particular Hamiltonian. These are three separate claims:
geometry, represented state-space action, and operator covariance. Each needs
its own hypotheses. Orbital labels, degeneracies, geometric stars, and calculator
symmetry flags do not supply electronic action matrices.

\begin{definition}[Finite group and composition convention]
\label{def:symmetry-finite-group}
A finite group $\mathcal G$ is a finite set of operations equipped with an
associative multiplication. Its identity element $e$ satisfies $eg=ge=g$,
and every element $g$ has an inverse $g^{-1}$ satisfying
$gg^{-1}=g^{-1}g=e$. Products remain in $\mathcal G$. The positive integer
$|\mathcal G|$ is the number of group elements, called the group order.
We read $gh$ as applying $h$ first and $g$ second. Thus a multiplication table
and its composition convention, not operation names alone, define the group
used in a calculation.
\end{definition}

\section{State spaces, representations and inner products}
\index{Gram matrix!state-space basis}

\begin{definition}[Source and target spaces of a state map]
\label{def:symmetry-source-target-spaces}
Choose two finite-dimensional complex inner-product state spaces
$\mathcal H_{\mathrm s}$ and $\mathcal H_{\mathrm t}$ in the declared model,
and a map $\widehat A:\mathcal H_{\mathrm s}\to\mathcal H_{\mathrm t}$.
The \emph{source space} is its domain $\mathcal H_{\mathrm s}$: the space
from which input states are taken. The \emph{target space} is its codomain
$\mathcal H_{\mathrm t}$: the declared space containing the output states.
The target need not equal the image, which is the set of outputs actually
reached by this map. The subscripts $\mathrm s$ and $\mathrm t$ name these
roles, not species, symmetry sectors or particular physical approximations.
\end{definition}

Source and target are relative to a specified map; they are not additional
kinds of quantum state space. For an action within one space,
$\mathcal H_{\mathrm s}=\mathcal H_{\mathrm t}$, the same space has both
roles. For symmetry sewing between reciprocal fibers, they can instead be
the state spaces at two symmetry-related momenta. Neither their dimensions
nor the words source and target establish a physical identification between
them. Such an identification requires a declared map and its qualifications.

\begin{definition}[Coefficient space of an ordered state basis]
\label{def:symmetry-coefficient-space}
For each role $a\in\{\mathrm s,\mathrm t\}$, let
$\mathcal B_a=(b_{a,1},\ldots,b_{a,n_a})$ be an ordered basis of
$\mathcal H_a$, where $n_a=\dim\mathcal H_a$ is its dimension. Its
\emph{coefficient space} is $\mathcal V_a=\mathbb C^{n_a}$, consisting of
complex column vectors $x_a=(x_{a,1},\ldots,x_{a,n_a})^{\mathsf T}$.
Here $\mathsf T$ means transpose; no conjugation is used to arrange the
entries into a column. The basis reconstruction map $\iota_a$ associates
each such coefficient vector with a state:
\begin{equation}
  \iota_a:\mathcal V_a\longrightarrow\mathcal H_a,
  \qquad \iota_a(x_a)=\sum_{i=1}^{n_a}x_{a,i}b_{a,i}.
  \label{eq:symmetry-basis-reconstruction}
\end{equation}
The index $i$ enumerates basis vectors, and $x_{a,i}$ is the complex
coefficient multiplying the $i$th vector. Linear independence and spanning
make $\iota_a$ a linear bijection. The vector $x_a$ is a coordinate
representation of a state, not a real-space position, a basis vector itself,
or a column extracted from an unspecified data matrix. The state basis need
not be orthogonal.
\end{definition}

Thus the \emph{source coefficient space} is
$\mathcal V_{\mathrm s}=\mathbb C^{n_{\mathrm s}}$, and the
\emph{target coefficient space} is
$\mathcal V_{\mathrm t}=\mathbb C^{n_{\mathrm t}}$. The state map
$\widehat A$ has the coordinate form
$\iota_{\mathrm t}^{-1}\widehat A\iota_{\mathrm s}$: reconstruct the input
state, apply the state map, then express the output in the target basis.
If this map is complex-linear, its coordinate form is multiplication by the
matrix $M$ introduced below. This explains both the direction of $M$ and why
its columns and rows refer to different bases. Inner products throughout this
chapter are conjugate-linear in the first argument and linear in the second.

\begin{definition}[Ordinary complex-linear representation]
\label{def:symmetry-ordinary-representation}
On a coefficient space $\mathcal V=\mathbb C^n$ constructed from an ordered
state basis as above, a representation assigns to each $g\in\mathcal G$
an invertible $n\times n$ complex matrix $D(g)$. Here $n$ is the dimension
of the represented space. Its action sends a coefficient vector
$x\in\mathcal V$ to $D(g)x$. Complex linearity means that sums and complex
scalar multiples of coefficient vectors are preserved. The assignment
respects group composition:
\begin{equation}
  D(e)=I,\qquad D(gh)=D(g)D(h).
  \label{eq:symmetry-ordinary-law}
\end{equation}
Here $I$ is the identity matrix on $\mathcal V$; the equations also imply
$D(g^{-1})=D(g)^{-1}$. The word \emph{ordinary} excludes an additional scalar
phase in the multiplication law. Distinct group elements may have the same
matrix: the representation is \emph{faithful} only when this never happens.
Faithfulness is not assumed in this chapter.
\end{definition}

A represented action includes both its space and its assignment of matrices.
The same geometric group can act differently on different state spaces.
For example, an action on atomic trial coordinates and one on retained band
coordinates are not identified merely because both have dimension ten.
Reciprocal fibers, reciprocal closure and translation phases require additional
structure, addressed in Chapter~\ref{ch:training-sewing-periodic-fibers}.

\begin{definition}[State-space basis Gram matrix]
\label{def:symmetry-state-gram}
The matrix $\mathsf G_a\in\mathbb C^{n_a\times n_a}$ has entries
$(\mathsf G_a)_{ij}=\langle b_{a,i},b_{a,j}\rangle_{\mathcal H_a}$.
Consequently the represented inner product and norm are
\begin{equation}
  \langle x,y\rangle_{\mathsf G_a}=x^{\dagger}\mathsf G_a y,
  \qquad \|x\|_{\mathsf G_a}=\sqrt{x^{\dagger}\mathsf G_a x}.
  \label{eq:symmetry-coordinate-inner-product}
\end{equation}
Here $x,y\in\mathcal V_a$ are coefficient columns, and $\dagger$ on a matrix
or column means conjugate transpose. A Gram matrix of a basis is Hermitian,
meaning $\mathsf G_a^{\dagger}=\mathsf G_a$, and positive definite, meaning
$x^{\dagger}\mathsf G_a x>0$ for every nonzero $x$. These properties ensure
that the squared norm is real and positive and that $\mathsf G_a$ is invertible.
An orthonormal basis has $\mathsf G_a=I$; a square array alone does not justify
this choice.
\end{definition}

Here basis Gram entries and coefficient columns are dimensionless. This
state-space metric is not the spatial lattice metric
$G_{ij}=\mathbf a_i\cdot\mathbf a_j$: the latter pairs direct-lattice vectors
$\mathbf a_i$ and has units of length squared.

\begin{definition}[Metric adjoint, isometry and unitary identification]
\label{def:symmetry-metric-map}
A complex-linear map from $\mathcal V_{\mathrm s}$ to $\mathcal V_{\mathrm t}$
has a coefficient matrix $M\in\mathbb C^{n_{\mathrm t}\times n_{\mathrm s}}$.
Its metric adjoint $M^{\sharp}$ maps in the reverse direction and is defined
by $\langle Mx,y\rangle_{\mathsf G_{\mathrm t}}
=\langle x,M^{\sharp}y\rangle_{\mathsf G_{\mathrm s}}$ for every coefficient
vector $x\in\mathcal V_{\mathrm s}$ representing a source state and every
coefficient vector $y\in\mathcal V_{\mathrm t}$ representing a target state. The map is an \emph{isometry} when it
preserves these inner products. In coordinates these statements give
\begin{equation}
  M^{\sharp}=\mathsf G_{\mathrm s}^{-1}M^{\dagger}\mathsf G_{\mathrm t},
  \qquad M^{\dagger}\mathsf G_{\mathrm t}M=\mathsf G_{\mathrm s}.
  \label{eq:symmetry-metric-isometry}
\end{equation}
Thus $M^{\sharp}$ has size $n_{\mathrm s}\times n_{\mathrm t}$.
An isometry is injective but need not fill the target. It is a \emph{unitary
identification} when it is also surjective; for equal finite source and target
dimensions, isometry already implies surjectivity. A metric-preserving
representation means that every $D(g)$ is unitary with respect to the stated
Gram matrix, not necessarily with respect to the Euclidean coefficient norm.
\end{definition}

Inverse notation here and below states identities, not a prescription for
explicit numerical inversion. To change source or target bases, let $A_a$ be
an invertible $n_a\times n_a$ matrix whose columns are the new basis vectors
in the old coordinates. Old and new columns are related by $x_a=A_a x'_a$.
For a complex-linear map, the same physical inner products and map are then
represented by
\begin{equation}
  \mathsf G'_a=A_a^{\dagger}\mathsf G_a A_a,
  \qquad M'=A_{\mathrm t}^{-1}MA_{\mathrm s}.
  \label{eq:symmetry-basis-change}
\end{equation}
Primes denote the new coordinates, not a different physical action. Changing
$M$ without changing the associated Gram matrices generally changes the problem.

\begin{definition}[Coefficient-action operator and covariance]
\label{def:symmetry-operator-covariance}
Let $h_a\in\mathbb C^{n_a\times n_a}$ represent a Hamiltonian acting on
$\mathcal H_a$: acting on a state with column $x$ produces a state with
column $h_a x$. The operator is \emph{self-adjoint} when its metric adjoint
equals itself. The source and target Hamiltonians are \emph{covariant under
$M$} when mapping and Hamiltonian evolution agree in either order.
In coordinates, source self-adjointness and covariance mean, respectively,
\begin{equation}
  h_{\mathrm s}^{\dagger}\mathsf G_{\mathrm s}
    =\mathsf G_{\mathrm s}h_{\mathrm s},
  \qquad h_{\mathrm t}M=Mh_{\mathrm s}.
  \label{eq:symmetry-operator-covariance}
\end{equation}
Target self-adjointness has the analogous equation with subscript $\mathrm t$.
When source and target are the same space and Hamiltonian, covariance reduces
to commutation, $hM=Mh$. In that case a qualified unitary $M$ is a symmetry
of that Hamiltonian, not merely an inner-product-preserving change of states.
Compatible physical conventions, including energy references, remain part of
the comparison.
\end{definition}

The matrix element $\langle b_{a,i},\widehat H_a b_{a,j}\rangle$ is instead
entry $(i,j)$ of $\mathsf G_a h_a$, where $\widehat H_a$ denotes the abstract
Hamiltonian. This matrix is Hermitian for a self-adjoint Hamiltonian, whereas
$h_a$ need not be Euclidean Hermitian. Both carry energy units. Confusing
matrix elements with coefficient actions changes the covariance equation.

\begin{definition}[Antiunitary identification]
\label{def:symmetry-antiunitary-map}
Let $\mathcal A:\mathcal V_{\mathrm s}\to\mathcal V_{\mathrm t}$ be the
coordinate form of a conjugate-linear state map. For any two coefficient
vectors $x,y\in\mathcal V_{\mathrm s}$ representing source states and any
complex scalars $\alpha,\beta$, conjugate linearity means
$\mathcal A(\alpha x+\beta y)=\overline\alpha\,\mathcal A x
+\overline\beta\,\mathcal A y$.
Unlike $y$ in the adjoint identity above, both vectors here belong to the
source coefficient space. An overline means complex conjugation. In coordinates write
$\mathcal A=MK$, so that $\mathcal A x=M\overline x$, where
$Kx=\overline x$ is componentwise conjugation in the declared source basis
and $M$ is the $n_{\mathrm t}\times n_{\mathrm s}$ matrix multiplying the
conjugated column. The map is \emph{antiunitary} when it is bijective
and norm-preserving; without surjectivity it is only a conjugate-linear
isometry. Norm preservation requires
\begin{equation}
  M^{\dagger}\mathsf G_{\mathrm t}M=\overline{\mathsf G_{\mathrm s}}.
  \label{eq:symmetry-antiunitary-metric}
\end{equation}
Indeed the transformed squared norm is
$x^{\mathsf T}M^{\dagger}\mathsf G_{\mathrm t}M\overline{x}$, which equals
$x^{\dagger}\mathsf G_{\mathrm s}x$ under this condition. Antiunitary operator
covariance uses $h_{\mathrm t}M=M\overline{h_{\mathrm s}}$, not the second
identity in Eq.~\eqref{eq:symmetry-operator-covariance}. Products must retain
conjugation: $(MK)(NK)=M\overline N$. Antiunitary sector theory is outside the
ordinary representation results below. In the product above $N$ is the
coefficient matrix of a second, composable conjugate-linear map, and each
conjugation acts in its corresponding coordinate space.
\end{definition}

\begin{lemma}[Change of coordinates for a conjugate-linear map]
\label{lem:symmetry-antiunitary-basis-change}
Let $A_{\mathrm s}$ and $A_{\mathrm t}$ be the invertible source and target
basis matrices defined in Eq.~\eqref{eq:symmetry-basis-change}. If the old
coordinate action is $x\mapsto M\overline x$, its new matrix is
\begin{equation}
  M'=A_{\mathrm t}^{-1}M\overline{A_{\mathrm s}}.
  \label{eq:symmetry-antiunitary-basis-change}
\end{equation}
Here $M'$ multiplies the conjugate of the new source coefficient vector;
the transformed Gram matrices remain
$\mathsf G'_a=A_a^{\dagger}\mathsf G_a A_a$.
\end{lemma}
\begin{proof}
For a new source vector $x'$, the old vector is $A_{\mathrm s}x'$.
Transforming its output to the new target coordinates gives
$A_{\mathrm t}^{-1}M\overline{A_{\mathrm s}x'}
=A_{\mathrm t}^{-1}M\overline{A_{\mathrm s}}\,\overline{x'}$,
which proves the formula. If the original map preserves norms, substitution
also gives
\begin{equation}
  (M')^{\dagger}\mathsf G'_{\mathrm t}M'
  =A_{\mathrm s}^{\mathsf T}\overline{\mathsf G_{\mathrm s}}
      \overline{A_{\mathrm s}}
  =\overline{\mathsf G'_{\mathrm s}}.
  \label{eq:symmetry-antiunitary-metric-basis-change}
\end{equation}
Thus the antiunitary metric condition is preserved by the same coordinate
change; surjectivity is also unchanged because the basis maps are bijective.
\end{proof}

\paragraph{Mechanization boundary.}
External Lean evidence for the coordinate basis-change identity, algebraic
metric-equation preservation and conjugate-linear composition is described in
Section~\ref{sec:mechanization-conjugate-linear-coordinate-algebra}, including
its exact hypotheses and provisional source bindings. This does not mechanize
the norm or physical interpretation, the sector results, or the whole chapter.

For an analytical check, start with componentwise conjugation on $\mathbb C^2$,
so $M=I_2$, and change both bases by
$A_{\mathrm s}=A_{\mathrm t}=\operatorname{diag}(i,1)$, where $i^2=-1$.
The correct new matrix is $\operatorname{diag}(-1,1)$; incorrectly using the
linear formula gives $I_2$. Both preserve the identity metric, so a norm
check alone would not detect the wrong action. For the new input vector
$(1,1)^{\mathsf T}$, their outputs are respectively $(-1,1)^{\mathsf T}$
and $(1,1)^{\mathsf T}$, which represent different states even up to a common
nonzero scalar. This illustrates why the conjugation convention belongs to
the action's definition, not just to a numerical norm test.

\section{Closure precedes an induced representation}
\index{subspace!symmetry closure}

\begin{definition}[Retained frame and induced metric]
\label{def:symmetry-retained-frame}
For each source or target space, choose a positive retained dimension
$r_a\leq n_a$. A \emph{frame} here is an $n_a\times r_a$ matrix $F_a$ whose
columns are linearly independent ambient coefficient columns; it is not a
redundant frame. The notation $\operatorname{Ran}F_a$ means the set of all
columns $F_a u$ with $u\in\mathbb C^{r_a}$. This range is the retained
subspace $\mathcal S_a\subseteq\mathcal V_a$, whereas $u$ is its coordinate
column in the frame. The induced Gram matrix $\mathsf G_{F_a}$ has size
$r_a\times r_a$ and pairs retained coordinate columns so that their inner
product equals that of the ambient columns they represent.
\end{definition}

\begin{definition}[Metric projection, compressed map and leakage]
\label{def:symmetry-compression-leakage}
The projector $P_{F_a}$ maps ambient columns onto $\mathcal S_a$, orthogonally
in the ambient metric $\mathsf G_a$. For an ambient map $M$, the compressed
matrix $D$ maps source retained coordinates to the target coordinates of the
metric projection of $MF_{\mathrm s}u$. The residual matrix $L$, called
\emph{leakage}, maps the same source coordinates to the portion left outside
the target subspace. Their formulas are
\begin{align}
  \mathsf G_{F_a}&=F_a^{\dagger}\mathsf G_a F_a,
    &a&\in\{\mathrm s,\mathrm t\},
    \label{eq:symmetry-frame-metric}\\
  P_{F_a}&=F_a\mathsf G_{F_a}^{-1}F_a^{\dagger}\mathsf G_a,
    \label{eq:symmetry-frame-projector}\\
  D&=\mathsf G_{F_{\mathrm t}}^{-1}
      F_{\mathrm t}^{\dagger}\mathsf G_{\mathrm t}MF_{\mathrm s},
    \label{eq:symmetry-compressed-action}\\
  L&=MF_{\mathrm s}-F_{\mathrm t}D.
    \label{eq:symmetry-leakage}
\end{align}
In these formulas $P_{F_a}$ has size $n_a\times n_a$, $D$ has size
$r_{\mathrm t}\times r_{\mathrm s}$, and $L$ has size
$n_{\mathrm t}\times r_{\mathrm s}$. Full column rank and a positive-definite
ambient metric make $\mathsf G_{F_a}$ positive definite. The projector
satisfies $P_{F_a}^2=P_{F_a}$, has range $\mathcal S_a$, and is self-adjoint
in that metric. It need not be a Euclidean Hermitian matrix.
For each source column $u$, $Du$ is the unique target column minimizing
$\|MF_{\mathrm s}u-F_{\mathrm t}v\|_{\mathsf G_{\mathrm t}}$ over target
retained columns $v$. Orthogonality of the minimizing residual gives
$F_{\mathrm t}^{\dagger}\mathsf G_{\mathrm t}L=0$.
The symbol $D$ here denotes a compressed map, not yet a representation.
\end{definition}

\begin{definition}[Closure and invariant subspace]
\label{def:symmetry-subspace-closure}
The map \emph{closes from the source to the target retained subspace} when
$M\mathcal S_{\mathrm s}\subseteq\mathcal S_{\mathrm t}$, equivalently
$L=0$. Equality is a stronger claim unless dimensions and injectivity supply
it. For an action on one ambient space, a subspace $\mathcal S$ is
\emph{invariant} under the group when $M(g)\mathcal S\subseteq\mathcal S$
for every $g\in\mathcal G$. Since inverses are also included, group
invariance implies $M(g)\mathcal S=\mathcal S$. This is invariance of a
space, not a requirement that each of its vectors be fixed.
\end{definition}

For an antiunitary action, the source term in the leakage is instead
$M\overline{F_{\mathrm s}}$, with an accompanying conjugate-linear retained map.

\begin{proposition}[Metric closure criterion]
If $M$ is an isometry and the frames have equal finite rank, the compressed
linear map is an isometry exactly when $M\operatorname{Ran}F_{\mathrm s}
=\operatorname{Ran}F_{\mathrm t}$.
\end{proposition}
\begin{proof}
Orthogonality of the leakage gives
\begin{equation}
  D^{\dagger}\mathsf G_{F_{\mathrm t}}D
  =\mathsf G_{F_{\mathrm s}}-L^{\dagger}\mathsf G_{\mathrm t}L.
  \label{eq:symmetry-leakage-defect}
\end{equation}
Positive definiteness makes the last term zero exactly when $L=0$.
This gives inclusion of the transformed source subspace in the target;
isometry and equal rank promote inclusion to equality.
\end{proof}

\begin{definition}[Restriction versus compression]
\label{def:symmetry-restriction-compression}
On one ambient space, the \emph{restriction} $M|_{\mathcal S}$ applies $M$
only to source vectors in a chosen subspace $\mathcal S$. Its outputs still
belong to the ambient space unless $\mathcal S$ is invariant. In contrast,
\emph{compression} onto $\mathcal S$ applies the metric projector
$P_{\mathcal S}$ after $M$, giving $P_{\mathcal S}M|_{\mathcal S}$.
This always maps into $\mathcal S$, but can discard part of the transformed
state. Restriction and compression represent the same map on $\mathcal S$
exactly when $M\mathcal S\subseteq\mathcal S$.
\end{definition}

To obtain a retained representation, let $M(g)$ be an ordinary ambient
representation and let one frame $F$ span a subspace invariant under every
$M(g)$. Closure gives $M(g)F=FD(g)$. For group elements $g,h$, acting twice
on a retained column gives
\begin{equation}
  FD(gh)=M(gh)F=M(g)M(h)F=FD(g)D(h).
  \label{eq:symmetry-restricted-group-law}
\end{equation}
The independent columns of $F$ make the coordinate map injective, so one can
cancel $F$ and obtain the group law for $D$. The same argument at the identity
gives $D(e)=I_r$, where $r$ is the number of columns of $F$. Thus the
compressed matrices become the coordinates of a genuine restricted
representation only after closure is established. A polar factor or
pseudoinverse does not establish the missing invariance.

For a metric-unitary ambient representation, invariance also has an equivalent
projector formulation:
\begin{equation}
  M(g)P_{\mathcal S}=P_{\mathcal S}M(g)
  \quad\hbox{for every }g\in\mathcal G.
  \label{eq:symmetry-invariant-projector-commutation}
\end{equation}
To see why, let $x$ be orthogonal to $\mathcal S$ and let $y\in\mathcal S$.
Unitarity gives $\langle M(g)x,y\rangle
=\langle x,M(g)^{-1}y\rangle=0$, because inverse group actions also preserve
$\mathcal S$. Hence both the subspace and its metric-orthogonal complement
are invariant, so projection commutes with the action. Conversely,
commutation gives $M(g)P_{\mathcal S}x=P_{\mathcal S}M(g)x\in\mathcal S$,
which proves invariance. The inner products in this argument are the same
ambient metric inner products, not Euclidean products introduced by convenience.

\section{Overlap forms and intertwiners}
\index{intertwiner|textbf}

\begin{definition}[Intertwiner of represented actions]
\label{def:symmetry-intertwiner}
Let $\mathcal C=\mathbb C^{r_C}$ and $\mathcal W=\mathbb C^{r_W}$ be identified
coefficient spaces carrying representations $D_C(g)$ and $D_W(g)$ of the
same group. An \emph{intertwiner} is a complex-linear map
$T:\mathcal C\to\mathcal W$ satisfying $D_W(g)T=TD_C(g)$ for every group
element. Its $r_W\times r_C$ matrix may be rectangular, noninvertible or zero.
The condition means that applying the source action before $T$ agrees with
applying the target action afterwards. It asserts neither isometry nor
Hamiltonian covariance, which involve additional structures.
\end{definition}

For a construction that may yield such a map, use one ambient coefficient
space $\mathbb C^n$ with positive-definite metric $\mathsf G$. Let
$W\in\mathbb C^{n\times r_W}$ and $C\in\mathbb C^{n\times r_C}$ be
full-column-rank frames. Here $C$ and $W$ are frame matrices, while
$\mathcal C$ and $\mathcal W$ are their retained coordinate spaces, not
additional ambient spaces. Overlaps and orthogonal projections can be defined
without a symmetry action and without an invariance assumption. Whether the
resulting map intertwines qualified actions is a separate question.

\begin{definition}[Overlap form and projected coefficient map]
\label{def:symmetry-overlap-form}
The overlap matrix $B\in\mathbb C^{r_W\times r_C}$ and projected coefficient
map $T$ are
\begin{equation}
  B=W^{\dagger}\mathsf G C,\qquad
  T=\mathsf G_W^{-1}B,\qquad \mathsf G_W=W^{\dagger}\mathsf G W.
  \label{eq:symmetry-overlap-coefficient-map}
\end{equation}
For columns $u\in\mathcal W$ and $v\in\mathcal C$, the scalar
$u^{\dagger}Bv$ is the ambient inner product $\langle Wu,Cv\rangle_{\mathsf G}$.
Thus $B$ stores pairings, rather than the coordinates of projected states.
The frame metric $\mathsf G_W$ is $r_W\times r_W$; solving
$\mathsf G_W Tv=Bv$ converts the overlap pairings to coefficients. In
particular, $WTv=P_W Cv$, where $P_W$ is the metric projector built from $W$.
If $\mathsf G_W=I$, then $T=B$; otherwise there is no general identification
of the two, although they can coincide for particular overlaps.
\end{definition}

\begin{proposition}[Overlap invariance and coefficient intertwining]
Suppose additionally that both frame ranges are invariant under the same
metric-preserving ordinary ambient representation $M(g)$. Define their
induced representations by $M(g)C=CD_C(g)$ and $M(g)W=WD_W(g)$. Then
\begin{equation}
  D_W(g)^{\dagger}BD_C(g)=B,
  \qquad D_W(g)T=TD_C(g).
  \label{eq:symmetry-overlap-intertwiner}
\end{equation}
\end{proposition}
\begin{proof}
Insert $M(g)W=WD_W(g)$ and $M(g)C=CD_C(g)$ in
$W^{\dagger}M(g)^{\dagger}\mathsf G M(g)C=B$ to obtain the first identity.
Metric preservation on $W$ gives
$D_W(g)^{\dagger}\mathsf G_W=\mathsf G_WD_W(g)^{-1}$.
Substitution of $B=\mathsf G_WT$ yields the second identity. The same algebra
shows their equivalence for metric-preserving $D_W(g)$.
\end{proof}

\subsection{Normalized overlaps and principal angles}

The distinction between pairings and coordinate maps also matters for angles.
Let $\mathsf G_C=C^{\dagger}\mathsf G C$ be the positive-definite induced
metric of the $C$ frame, complementing $\mathsf G_W$ defined above. For either
frame metric, the power $-1/2$ denotes the inverse of its unique positive-definite
Hermitian square root. The normalized overlap is
\begin{equation}
  \widetilde B=\mathsf G_W^{-1/2}B\mathsf G_C^{-1/2}.
  \label{eq:symmetry-normalized-overlap}
\end{equation}
It pairs the metric-orthonormal frames $W\mathsf G_W^{-1/2}$ and
$C\mathsf G_C^{-1/2}$ spanning the original subspaces. It does not change
those subspaces or select a new retained space.

\begin{definition}[Principal angles between represented subspaces]
\label{def:symmetry-principal-angles}
Write $p=\min(r_W,r_C)$ for the smaller frame dimension and let
$s_1\geq\cdots\geq s_p$ be the singular values of $\widetilde B$.
The \emph{principal angles} are the numbers
$0\leq\theta_1\leq\cdots\leq\theta_p\leq\pi/2$ defined by
$\cos\theta_i=s_i$. Their cosines lie in $[0,1]$ because the frames are
orthonormal in the same ambient inner product. Angles are unchanged by
invertible changes of coordinates within either frame's range, provided the
metrics are transformed consistently. Equal-dimensional subspaces have all
principal angles zero exactly when the subspaces coincide. With unequal
dimensions, all zero angles establish containment of the smaller subspace,
not equality of the two spaces.
\end{definition}

For a simple analytical illustration, use a one-dimensional ambient space
with metric $\mathsf G=(1)$, and choose frames $W=(2)$ and $C=(1)$.
Both span the whole space, but $\mathsf G_W=(4)$, $\mathsf G_C=(1)$,
$B=(2)$, $T=(1/2)$ and $\widetilde B=(1)$. The raw overlap singular value
is two, whereas the sole principal angle is zero. The projected coefficient
$1/2$ is also correct: multiplying it by the frame vector two reproduces the
state with ambient coefficient one. Thus raw overlap values, projection
coefficients and angle cosines answer different questions. Normalization is
an explicit mathematical conversion, not permission to apply an unchanged
raw-coordinate gate to a different numerical quantity.

\section{Isotypic sectors and a conditional rank obstruction}
\index{isotypic projector|textbf}

\begin{definition}[Irreducibility and equivalence]
\label{def:symmetry-irreducibility}
A nonzero represented space is \emph{irreducible} when its only invariant
subspaces are the zero space and the whole space. An \emph{irrep} is an
irreducible representation. Two representations are \emph{equivalent} when
an invertible intertwiner identifies them. Equivalence is a statement about
all group actions under one consistent change of basis, not merely equal
dimension or a shared eigenvalue of a Hamiltonian.
\end{definition}

\begin{lemma}[Invariant inner product and complete reducibility]
\label{lem:symmetry-complete-reducibility}
Every ordinary finite-dimensional complex representation of a finite group
admits an invariant positive-definite inner product. It is \emph{completely
reducible}: its space is a direct sum of invariant irreducible subspaces.
\end{lemma}
\begin{proof}
Let $D(g)$ act on a coefficient space $\mathcal V$, and choose any
positive-definite inner product $\langle\cdot,\cdot\rangle_0$ on that space.
For vectors $x,y\in\mathcal V$, define its group average by
\begin{equation}
  \langle x,y\rangle_{\mathrm{av}}
    =\frac{1}{|\mathcal G|}\sum_{g\in\mathcal G}
         \langle D(g)x,D(g)y\rangle_0.
  \label{eq:symmetry-averaged-inner-product}
\end{equation}
The subscript $\mathrm{av}$ denotes this averaged inner product, not a
statistical expectation. A sum of positive squared norms is positive for
nonzero $x$, so the average is positive definite. For each group element
$h$, the bijection $g\mapsto gh$ of $\mathcal G$ proves
$\langle D(h)x,D(h)y\rangle_{\mathrm{av}}
=\langle x,y\rangle_{\mathrm{av}}$.
An invariant subspace therefore has an invariant orthogonal complement, by
the argument in Eq.~\eqref{eq:symmetry-invariant-projector-commutation}.
Splitting a nontrivial invariant subspace from its complement and inducting
on dimension yields a direct sum of irreducibles; the zero-dimensional space
is the empty sum. This construction proves existence of an invariant metric.
It does not replace the specified physical metric in a numerical check.
\end{proof}

\begin{lemma}[Intertwiners between irreducible representations]
\label{lem:symmetry-schur}
A nonzero intertwiner between two finite-dimensional complex irreducible
representations is an isomorphism. An intertwining endomorphism of one such
representation is a scalar multiple of the identity.
\end{lemma}
\begin{proof}
The kernel and image of an intertwiner are invariant. Irreducibility and a
nonzero map make the kernel zero and the image the whole target, proving
invertibility. For an intertwining endomorphism $T$ on a nonzero complex
finite-dimensional space, choose an eigenvalue $\lambda\in\mathbb C$.
The nonzero eigenspace $\ker(T-\lambda I)$ is invariant and hence is the
whole space. Thus $T=\lambda I$. In particular, the vector space of
intertwiners between inequivalent irreducibles has dimension zero; between
a chosen irreducible and itself it has dimension one.
\end{proof}

\begin{definition}[Character and irreducible multiplicity]
\label{def:symmetry-character-multiplicity}
For an ordinary representation $D(g)$ on a retained coefficient space
$\mathbb C^r$, its \emph{character} is $\chi_D(g)=\operatorname{tr}D(g)$,
where the trace is the sum of diagonal entries. Trace is invariant under a
change of basis, so the character does not depend on a chosen frame within
that represented space. Choose an index set $\mathcal I$ containing one
representative of every equivalence class of complex irreps of $\mathcal G$.
For $\alpha\in\mathcal I$, let $V_\alpha$ be its irreducible coefficient
space, $d_\alpha=\dim V_\alpha$ its positive dimension, and
$D_\alpha(g)$ its action matrix. Its character is
$\chi_\alpha(g)=\operatorname{tr}D_\alpha(g)$.

The \emph{multiplicity} $m_{F,\alpha}$ is the number of copies of
$V_\alpha$ occurring in the representation on the coordinates of frame $F$.
Lemma~\ref{lem:symmetry-complete-reducibility} gives a change of coordinates
with
\begin{equation}
  \mathbb C^r\simeq\bigoplus_{\alpha\in\mathcal I}
     \left(V_\alpha\otimes\mathbb C^{m_{F,\alpha}}\right),
  \qquad r=\sum_{\alpha\in\mathcal I}d_\alpha m_{F,\alpha}.
  \label{eq:symmetry-isotypic-decomposition}
\end{equation}
Here $\simeq$ means representation equivalence, $\bigoplus$ is a direct sum,
and $\otimes$ is a tensor product. The group acts on $V_\alpha$ and as the
identity on the multiplicity space $\mathbb C^{m_{F,\alpha}}$; a zero
multiplicity means that irrep is absent. In this notation $F$ identifies the
frame carrying $D$, not a group element.
\end{definition}

\begin{lemma}[Orthogonality of irreducible characters]
\label{lem:symmetry-character-orthogonality}
For complex-valued functions $u,v:\mathcal G\to\mathbb C$, define the
normalized character-space inner product by
\begin{equation}
  \langle u,v\rangle_{\mathcal G}
    =\frac{1}{|\mathcal G|}\sum_{g\in\mathcal G}\overline{u(g)}v(g).
  \label{eq:symmetry-character-inner-product}
\end{equation}
This is an inner product on functions of group elements, not a state-space
Gram metric. For the chosen inequivalent irreps it obeys
$\langle\chi_\alpha,\chi_\beta\rangle_{\mathcal G}=\delta_{\alpha\beta}$,
where $\delta_{\alpha\beta}$ is one when $\alpha=\beta$ and zero otherwise.
\end{lemma}
\begin{proof}
Let $\operatorname{Hom}(V_\alpha,V_\beta)$ denote the vector space of all
complex-linear maps between the two irreducible spaces. Its elements have
$d_\beta\times d_\alpha$ matrices $X$. Define a linear averaging operator
$Q_{\alpha\beta}$ on this space by
\begin{equation}
  Q_{\alpha\beta}(X)=\frac{1}{|\mathcal G|}\sum_{g\in\mathcal G}
      D_\beta(g)X D_\alpha(g)^{-1}.
  \label{eq:symmetry-intertwiner-average}
\end{equation}
Reindexing by left multiplication of group elements shows that its output
intertwines the two representations. Conversely, every intertwiner is fixed
by this average. Hence $Q_{\alpha\beta}$ is an idempotent operator whose
image is precisely the intertwiner space. By Lemma~\ref{lem:symmetry-schur},
that image has dimension $\delta_{\alpha\beta}$. An idempotent is the
identity on its image and zero on its kernel, so its trace equals this
dimension.

For square matrices $A$ on $V_\beta$ and $B$ on $V_\alpha$, the operator
$X\mapsto AXB$ on this matrix space has trace
$(\operatorname{tr}A)(\operatorname{tr}B)$: in the basis of matrices with
one entry equal to one, its diagonal coefficient at position $(i,j)$ is
$A_{ii}B_{jj}$, and summing over $i,j$ gives the product. Applying this fact
to the average computes the same trace as
$|\mathcal G|^{-1}\sum_g\chi_\beta(g)\operatorname{tr}(D_\alpha(g)^{-1})$.
Using the invariant inner product of
Lemma~\ref{lem:symmetry-complete-reducibility}, the irreducible matrices may
be written in unitary coordinates. Their inverse
traces are therefore $\overline{\chi_\alpha(g)}$; trace is unchanged by
returning to the original coordinates. Equating the two trace expressions
proves the asserted orthogonality.
\end{proof}

Orthogonality also shows that there are at most $|\mathcal G|$ inequivalent
irreps: their nonzero, mutually orthogonal characters belong to the
$|\mathcal G|$-dimensional space of complex-valued functions on the group.
Thus the complete irrep index set $\mathcal I$ is finite.

\begin{definition}[Isotypic sector]
\label{def:symmetry-isotypic-sector}
The \emph{isotypic sector} for $\alpha$ is the sum of all invariant copies
equivalent to $V_\alpha$; its dimension is $d_\alpha m_{F,\alpha}$, not just
the number of copies. This sector does not select a unique basis or a unique
splitting of repeated copies.
\end{definition}

\begin{proposition}[Character projector onto an isotypic sector]
\label{prop:symmetry-character-projector}
For an ordinary representation $D(g)$ preserving the induced frame metric
$\mathsf G_F$, the metric-orthogonal projector onto sector $\alpha$ is
\begin{equation}
  \Pi_\alpha=\frac{d_\alpha}{|\mathcal G|}
     \sum_{g\in\mathcal G}\overline{\chi_\alpha(g)}D(g).
  \label{eq:symmetry-isotypic-projector}
\end{equation}
The sum is over every group element, and the overline conjugates the
character, not the representation matrix. Each $\Pi_\alpha$ has size
$r\times r$. It is the identity on copies of irrep $\alpha$ and zero on
inequivalent irreps. Consequently
$\Pi_\alpha^2=\Pi_\alpha$ (idempotence),
$\Pi_\alpha\Pi_\beta=0$ for distinct labels $\alpha,\beta$ (mutual
orthogonality), and $\sum_{\alpha\in\mathcal I}\Pi_\alpha=I_r$
(completeness). Orthogonality is in the frame metric:
$\Pi_\alpha^{\dagger}\mathsf G_F=\mathsf G_F\Pi_\alpha$.
\end{proposition}
\begin{proof}
Characters are invariant under conjugating group elements, because
$D_\alpha(hgh^{-1})=D_\alpha(h)D_\alpha(g)D_\alpha(h)^{-1}$ has the
same trace as $D_\alpha(g)$. Reindexing the projector sum therefore shows
that $\Pi_\alpha$ commutes with every group action. It also preserves each
invariant irreducible summand, since each term in the sum does.
On a copy of irrep $\beta$, Lemma~\ref{lem:symmetry-schur} makes its
restriction a scalar $c_{\alpha\beta}$ times the identity. Taking the trace
on that copy and applying Lemma~\ref{lem:symmetry-character-orthogonality}
gives $d_\beta c_{\alpha\beta}=d_\alpha\delta_{\alpha\beta}$.
Thus the scalar is one for $\beta=\alpha$ and zero otherwise, proving
idempotence, product orthogonality and completeness on the irreducible
decomposition.

Different isotypic sectors are orthogonal in the given invariant metric.
Indeed, metric projection onto an invariant irreducible copy commutes with
the group action by Eq.~\eqref{eq:symmetry-invariant-projector-commutation}.
Its restriction to an inequivalent irreducible copy is an intertwiner and
hence zero by Lemma~\ref{lem:symmetry-schur}. Summing such copies establishes
orthogonality of the sectors. Since $\Pi_\alpha$ acts as identity on one
sector and zero on all others, it is the metric-orthogonal projector and
satisfies the asserted metric self-adjointness identity.
\end{proof}

For exact ordinary representations with the qualified complete irrep list,
Lemma~\ref{lem:symmetry-character-orthogonality} also recovers each multiplicity:
\begin{equation}
  m_{F,\alpha}=\frac{1}{|\mathcal G|}
    \sum_{g\in\mathcal G}\overline{\chi_\alpha(g)}\chi_D(g).
  \label{eq:symmetry-character-multiplicity}
\end{equation}
To obtain this formula, take the trace of the direct-sum representation in
Eq.~\eqref{eq:symmetry-isotypic-decomposition}, giving
$\chi_D(g)=\sum_{\beta\in\mathcal I}m_{F,\beta}\chi_\beta(g)$, and pair
it with $\chi_\alpha$ using Eq.~\eqref{eq:symmetry-character-inner-product}.
Orthogonality removes every term except $m_{F,\alpha}$. Since the character
is independent of the decomposition, this formula also proves that the
multiplicities are unique even though the individual equivalent copies need
not be. This value is a nonnegative integer under those hypotheses. Supplied character
data, operation correspondence and the group law must first be qualified.
Rounding a computed trace or multiplicity to an integer does not establish
those hypotheses. An energy degeneracy likewise does not determine an irrep
or its multiplicity.

\begin{proposition}[Sector multiplicities bound intertwiner rank]
\label{prop:symmetry-sector-rank-bound}
Let $T:\mathcal C\to\mathcal W$ intertwine ordinary finite-dimensional
complex representations on identified coefficient spaces $\mathcal C$ and
$\mathcal W$ of a finite group. Write $m_{C,\alpha}$ and $m_{W,\alpha}$ for
the numbers of copies of the same irrep $V_\alpha$ in source and target.
The rank of $T$ is the dimension of its image; its kernel $\ker T$ is the
source subspace sent to zero. Both sums below run over the complete irrep
index set $\mathcal I$ introduced above. Then
\begin{align}
  \operatorname{rank}T
    &\leq\sum_\alpha d_\alpha
        \min(m_{C,\alpha},m_{W,\alpha}),
    \label{eq:symmetry-intertwiner-rank-bound}\\
  \dim\ker T
    &\geq\sum_\alpha d_\alpha
        \max(m_{C,\alpha}-m_{W,\alpha},0).
    \label{eq:symmetry-intertwiner-kernel-bound}
\end{align}
\end{proposition}
\begin{proof}
Decompose the source and target using
Lemma~\ref{lem:symmetry-complete-reducibility}. By
Lemma~\ref{lem:symmetry-schur}, every map between inequivalent irreducible
copies is zero, and every map between fixed equivalent copies is scalar
after their representations are identified. Therefore in irreducible
coordinates $T$ has blocks
$I_{d_\alpha}\otimes K_\alpha$, where $I_{d_\alpha}$ is the identity on
$V_\alpha$ and $K_\alpha$ is an $m_{W,\alpha}\times m_{C,\alpha}$ matrix
between the multiplicity spaces. Each block has rank at most
$d_\alpha\min(m_{C,\alpha},m_{W,\alpha})$. Summing and subtracting from
$\dim\mathcal C=\sum_\alpha d_\alpha m_{C,\alpha}$ proves both bounds.
\end{proof}

Matching multiplicities permit full-rank intertwiners but do not force them:
a block $K_\alpha$ may itself be deficient. For the overlap construction,
$\operatorname{rank}B=\operatorname{rank}T$ because $\mathsf G_W$ is invertible.
The proposition concerns exact ranks under exact representation and closure
hypotheses. Small numerical covariance residuals alone do not justify its
application as an exact rank certificate. Quantitative approximate versions
need separately derived perturbation bounds.

\subsection{A worked sector-mismatch example}

Consider a two-element group $\mathcal G=\{e,g\}$ with $g^2=e$.
Its two complex irreps are one-dimensional: label them $+$ and $-$, with
characters $\chi_+(e)=\chi_-(e)=1$ and $\chi_+(g)=1$, $\chi_-(g)=-1$.
Let the source and target coefficient spaces both be $\mathbb C^3$ with
identity metrics, but give them different actions:
\begin{equation}
  D_C(g)=\operatorname{diag}(1,1,-1),\qquad
  D_W(g)=\operatorname{diag}(1,-1,-1).
  \label{eq:symmetry-sector-example-actions}
\end{equation}
At $e$ both representations equal $I_3$. Their characters at $g$ are
$\chi_{D_C}(g)=1$ and $\chi_{D_W}(g)=-1$, while both identity characters
are three. Equation~\eqref{eq:symmetry-character-multiplicity} gives
$m_{C,+}=2$, $m_{C,-}=1$, $m_{W,+}=1$ and $m_{W,-}=2$.
For example, the source $+$ multiplicity is $(3+1)/2=2$, not the trace at
$g$ alone. The source isotypic projectors are
$(I_3+D_C(g))/2=\operatorname{diag}(1,1,0)$ for $+$ and
$(I_3-D_C(g))/2=\operatorname{diag}(0,0,1)$ for $-$.
The first selects a two-dimensional sector containing two equivalent copies;
it does not select a preferred one-dimensional basis within that sector.

The intertwining equation forces every source-to-target matrix to have the
form
\begin{equation}
  T=\begin{pmatrix}a&b&0\\0&0&c\\0&0&d\end{pmatrix},
  \qquad a,b,c,d\in\mathbb C.
  \label{eq:symmetry-sector-example-intertwiner}
\end{equation}
Here $a,b$ map the two source $+$ coordinates into the single target $+$
coordinate, while $c,d$ map the single source $-$ coordinate into the two
target $-$ coordinates. Opposite sectors cannot couple through an
intertwiner. The rank bound is therefore at most
$\min(2,1)+\min(1,2)=2$, and at least one source direction lies in the kernel
despite equal total dimensions. This upper bound is attained when both pairs
$(a,b)$ and $(c,d)$ are nonzero.

For contrast, take both actions to be $\operatorname{diag}(1,-1)$ on
$\mathbb C^2$ and choose $T=\operatorname{diag}(1,0)$. This map intertwines
the identical representations but has rank one. Equal multiplicities have
removed the forced sector mismatch, not every possible cause of rank loss.
Consequently an observed rank loss does not by itself prove a sector mismatch.
These are exact analytical examples, not silicon sector assignments or retained
numerical results.

\section{Analytical controls and numerical interpretation}

Two illustrative controls expose mistakes that orthonormal positive cases miss.
Take the two-element group $\mathcal G=\{e,g\}$ with $g^2=e$. On the
orthonormal coefficient space $\mathbb C^2$, let $D(e)=I_2$ and let $D(g)$
be the reflection matrix $R=\operatorname{diag}(1,-1)$, which fixes the first
coordinate and reverses the second. Let $A$ have columns equal to the new
basis vectors in these old coordinates. The new metric $\mathsf G$ and action
$M$ follow from Eq.~\eqref{eq:symmetry-basis-change}:
\begin{equation}
  A=\begin{pmatrix}1&1\\0&1\end{pmatrix},\quad
  \mathsf G=A^{\dagger}A=\begin{pmatrix}1&1\\1&2\end{pmatrix},\quad
  M=A^{-1}RA=\begin{pmatrix}1&2\\0&-1\end{pmatrix}.
  \label{eq:symmetry-nonorthogonal-control}
\end{equation}
Direct multiplication gives $M^2=I$ and $M^{\dagger}\mathsf G M=\mathsf G$,
but $M^{\dagger}M\ne I$. Euclidean unitarity would reject a valid action.
Conversely, in the original orthonormal coordinates retain only
$f=(1,1)^{\mathsf T}/\sqrt2$. Then $f^{\dagger}Rf=0$ and $Rf$ is orthogonal
to $f$. Compression of an exact ambient representation has lost closure and
is not a one-dimensional group representation. These are analytical examples,
not completed software tests or calculated silicon results.

\begin{definition}[Metric-compatible operator norm]
\label{def:symmetry-metric-operator-norm}
For a linear matrix $X$ from a nonzero source space with metric
$\mathsf G_{\mathrm s}$ to a target with metric $\mathsf G_{\mathrm t}$,
define
\begin{equation}
  \|X\|_{\mathsf G_{\mathrm t}\leftarrow\mathsf G_{\mathrm s}}
  =\sup_{x\ne0}\frac{\|Xx\|_{\mathsf G_{\mathrm t}}}
                         {\|x\|_{\mathsf G_{\mathrm s}}}.
  \label{eq:symmetry-metric-operator-norm}
\end{equation}
The supremum is over nonzero source coefficient columns. It measures the
largest target norm obtained per unit source norm. In particular, leakage
$L$ uses the source retained metric $\mathsf G_{F_{\mathrm s}}$ and the
target ambient metric $\mathsf G_{\mathrm t}$, not two ambient metrics of
unrelated sizes. Hamiltonian covariance residuals carry energy units, whereas
action and leakage residuals in the present convention are dimensionless.
\end{definition}

\subsection{Residuals must retain their mathematical type}

To specify reusable diagnostics, let $D(g)$ now denote candidate
$r\times r$ action matrices on one coefficient space with metric
$\mathsf G_F$, not matrices already assumed to pass the representation law.
For group elements $g_1,g_2$, define the identity residual $E_{\mathrm{id}}$,
composition residual $E_{\mathrm{law}}$, and isometry residual
$E_{\mathrm{iso}}$ by
\begin{align}
  E_{\mathrm{id}}&=D(e)-I_r,
    \label{eq:symmetry-identity-residual}\\
  E_{\mathrm{law}}(g_1,g_2)
    &=D(g_1g_2)-D(g_1)D(g_2),
    \label{eq:symmetry-group-law-residual}\\
  E_{\mathrm{iso}}(g)&=D(g)^{\sharp}D(g)-I_r.
    \label{eq:symmetry-isometry-residual}
\end{align}
Every residual in these formulas is a dimensionless coefficient map from the
same retained space to itself. The metric adjoint uses $\mathsf G_F$ on both
sides, so their operator norms also use $\mathsf G_F$ as both source and
target metric. In contrast, the array
$D(g)^{\dagger}\mathsf G_FD(g)-\mathsf G_F$ is an inner-product form.
Left multiplication by $\mathsf G_F^{-1}$ converts that form into
$E_{\mathrm{iso}}(g)$. An array's shape alone does not distinguish these
mathematical types or choose their norm.

For a candidate ambient map $M$ between source and target Hamiltonian spaces,
the covariance residual is the coefficient map
\begin{equation}
  E_{\mathrm{cov}}=h_{\mathrm t}M-Mh_{\mathrm s}.
  \label{eq:symmetry-covariance-residual}
\end{equation}
It has size $n_{\mathrm t}\times n_{\mathrm s}$ and energy units.
Its operator norm uses the ambient source and target metrics. It is not the
same object as the dimensionless leakage $L$, which uses a retained source
metric, or the overlap invariance residual
$D_W(g)^{\dagger}BD_C(g)-B$, which is a pairing between retained spaces.
A single untyped norm test cannot replace these distinctions.

Candidate projectors on one represented space have separate residuals:
$\Pi_\alpha^2-\Pi_\alpha$ for idempotence,
$\Pi_\alpha\Pi_\beta$ for distinct-sector orthogonality,
$\sum_{\alpha\in\mathcal I}\Pi_\alpha-I_r$ for completeness, and
$\Pi_\alpha^{\sharp}-\Pi_\alpha$ for metric self-adjointness.
Here $\alpha,\beta\in\mathcal I$ are distinct irrep labels when testing
orthogonality. The metric adjoint again belongs to the represented space;
Euclidean Hermiticity is not a substitute in a nonorthogonal frame.

Numerical qualification needs declared norms, conditioning and error bounds
for each of these quantities. An exact identity proved under exact hypotheses
is different from a bounded residual meeting a declared numerical tolerance.
An interval containing zero does not prove an exact residual is zero; an
interval straddling a decision threshold leaves that decision inconclusive.
No rounded character, repaired action or removed direction supplies a missing
hypothesis of the sector-rank proposition.
Input printing/parsing, arithmetic, finite-representation, model
reduction and parent-physics errors remain distinct. There is no universal
symmetry tolerance, silent polar repair, rank reduction, or permission to
apply a raw-overlap threshold to normalized singular values. Principal-angle
cosines require metric-orthonormal frames (or their explicitly justified
normalization); raw trial-coordinate overlaps need not have that meaning.
Higher precision addresses an actual arithmetic ambiguity, not missing native
digits or unqualified source conventions.

\section{Scientific and implementation boundary}

The first proposed reusable software slice is ordinary finite groups and
complex-linear metric-preserving actions. Antiunitary corepresentations,
projective factors, nonsymmorphic little-group theory and spinful double-valued
actions require separate contracts; the ordinary character formula does not
cover them by relabeling an operation. General sewing remains a separate
state-space identification category.

For silicon, source-bound parent, trial and Wannier actions, their metrics,
phases, closure and compatible character data must precede a sector diagnosis.
Appendix~\ref{app:preliminary-silicon-work} reports a verified numerical rank
stop, not an application of Proposition~\ref{prop:symmetry-sector-rank-bound}.
Geometric symmetry or a nearly vanishing overlap does not establish its
hypotheses. Native wavefunction extraction, reusable mathematics, and campaign
source admission and scientific disposition have distinct owners. Proposed
package incubation or later migration is not evidence that extraction or
symmetry-sector software already exists.
